\section{Related Work}

\label{sec:related}
% ══════════════════════════════════════════════════════════════════════════════

\subsection{Citizen Science Platforms}

Major biodiversity citizen science platforms include:
\begin{itemize}
  \item \textbf{eBird} \citep{sullivan2014ebird}: Structured bird
    checklists with 800M+ observations. Strong protocol design but
    limited to birds.
  \item \textbf{iNaturalist} \citep{inaturalist2024}: Multi-taxon
    platform with AI identification. Broad taxonomic coverage but
    unstructured opportunistic data.
  \item \textbf{Zooniverse} \citep{simpson2014zooniverse}: Platform
    for classification tasks (e.g., galaxy classification, camera
    trap identification).
  \item \textbf{GBIF} \citep{gbif2024}: Global aggregator of
    biodiversity data from multiple sources including citizen science.
\end{itemize}

\subsection{Gamification in Citizen Science}

\citet{bowser2013gamifying} identify four gamification strategies:
competition (leaderboards), cooperation (team goals), exploration
(discovery quests), and collection (digital specimen albums).
\citet{eveleigh2014designing} demonstrate that intrinsic motivation
(learning, contribution to science) outperforms extrinsic rewards for
sustained engagement. \citet{preece2019citizen} survey
gamification across 70 citizen science projects, finding that narrative
elements and progressive skill development are most effective for
long-term retention.

\subsection{Data Quality in Citizen Science}

\citet{kelling2015can} demonstrate that structured protocols (complete
checklists, standardized effort) substantially improve data quality.
\citet{johnston2020analytical} develop analytical methods for
accounting for observer skill in eBird data. \citet{isaac2014statistics}
formalize occupancy models that jointly estimate species presence
and detection probability from citizen science data.

\subsection{AI-Assisted Identification}

AI identification assistants in citizen science include iNaturalist's
computer vision \citep{van2021benchmarking}, Merlin for bird audio
\citep{kahl2021birdnet}, and PlantNet for plant identification
\citep{goeau2022overview}. These systems improve identification accuracy
but introduce new biases: AI suggestions may anchor observer decisions,
and AI errors propagate through community validation
\citep{terry2020think}.


% ══════════════════════════════════════════════════════════════════════════════
